\documentclass[11pt,a4paper]{article}
\usepackage[margin=2.5cm]{geometry}
\usepackage{parskip}
\usepackage{titlesec}
\usepackage{enumitem}
\usepackage{hyperref}
\usepackage{xcolor}

\definecolor{stepgray}{RGB}{90,90,100}

\titleformat{\section}{\large\bfseries}{}{0em}{}[\titlerule]
\titleformat{\subsection}{\normalsize\bfseries}{}{0em}{}

\setlist[itemize]{leftmargin=1.5em, itemsep=3pt}

\begin{document}

\begin{center}
  {\LARGE \textbf{Experiments Overview}} \\[4pt]
  {\small HGT Legal Outcome Prediction --- Explainability \& Structural Analysis}
\end{center}

\vspace{4pt}

This document summarises every experiment and test implemented for the explanation
phase of the thesis, written in pipeline order --- each step either produces outputs
that feed the next step, or deepens the interpretation of a previous step. The goal
is to record methodology and rationale in enough detail to support the research paper.

The underlying model is a Heterogeneous Graph Transformer (HGT) trained on Indian
Supreme Court cases. The graph contains case nodes, statute/provision/precedent nodes,
judge/court/party/lawyer nodes, and structured case section nodes (facts, arguments,
preamble, etc.). The model predicts case outcomes (petitioner win or loss). All
experiments operate on a frozen, trained model --- nothing is retrained.

\bigskip

% -----------------------------------------------------------------------
\section{Step 1 \texorpdfstring{\textcolor{stepgray}{---}}{---} Typed Counterfactual Masking (Core Explainability)}

\textbf{What we do.}
For every test case we first extract its \emph{L-hop induced receptive field}: the
complete set of nodes within $L$ hops of the target case in the heterogeneous graph
(where $L$ equals the number of GNN layers, typically 2), together with all directed
edges among those nodes. This is the exact subgraph the model used to compute the
prediction; we reconstruct it locally so the model can be run on it in isolation.

Within this local subgraph we enumerate \emph{mask groups}. Two kinds exist:
\begin{itemize}
  \item \textbf{Node groups} --- every non-case node in the receptive field (individual
    statutes, provisions, precedents, case sections, judges, courts, parties, lawyers).
    Masking removes all edges incident to that node in the local subgraph.
  \item \textbf{Relation-type groups} --- every distinct edge type (e.g.\ all
    ``cites-statute'' edges, all ``cited-by-section'' edges). Masking removes every
    edge of that type from the local subgraph.
\end{itemize}

For each group we run one forward pass of the frozen model with those edges removed
and record: the change in predicted-class probability (\emph{delta}), the change in
the probability of the positive label, and whether the prediction flips to the
opposite class. Groups are ranked by absolute delta. We simultaneously run the
unmasked forward pass with attention recording enabled and attach the mean HGT
attention weight assigned to each group's edges --- giving a second, attention-derived
importance signal for comparison.

The top-$k$ groups per case, aggregate evidence-type importance across all test cases,
aggregate relation-type importance, and a leakage-sensitivity summary (average delta
for identity-type nodes) are all written to disk. For each masked evidence node we
also compute a \emph{support profile}: how many training cases share that node and
what their label distribution (wins vs.\ losses) is. This directly links each
explanation back to the training data that shaped the model.

\textbf{Why.}
Standard attribution methods (SHAP, LIME, saliency maps) are designed for
fixed-input models. A GNN operating on a typed heterogeneous graph has no fixed
input structure: the number of evidence nodes varies per case, edges carry semantic
type, and message passing is inherently structural. Counterfactual masking is the
natural adaptation: we intervene directly on the graph structure and observe the
causal effect on the prediction. Masking at the \emph{edge} level (rather than
deleting nodes) ensures we capture the total contribution of each evidence node
regardless of which path through the graph it uses.

\textbf{What we can infer.}
\begin{itemize}
  \item \textbf{Per-case explanation:} Which specific statutes, provisions, or
    precedents drove each prediction. A positive delta means the group supported the
    original prediction; a negative delta means its presence worked against it.
  \item \textbf{Prediction flips:} When masking a single group flips the prediction
    entirely, that group is load-bearing. Flip rates across the test set reveal how
    often the model's decision depends on a single piece of evidence.
  \item \textbf{Relation-type importance:} The aggregate ranking of edge types shows
    which kinds of connections the model relies on most globally.
  \item \textbf{Attention--counterfactual agreement:} When both signals rank the same
    group highly, that is convergent evidence of importance. Disagreement raises
    questions about what HGT attention is actually tracking.
  \item \textbf{Early identity signal:} If judges, courts, petitioners, or lawyers
    appear persistently in top-ranked groups, the model may be using identity
    shortcuts. This is the seed for the shortcut and mask sensitivity experiments
    (Steps 9 and 10).
\end{itemize}

\textbf{Outputs that feed later steps:} top explanations CSV, case summary CSV,
counterfactual groups CSV, connected-case label distribution CSV, leakage sensitivity
summary CSV, attention--counterfactual overlap CSV.

\bigskip

% -----------------------------------------------------------------------
\section{Step 2 \texorpdfstring{\textcolor{stepgray}{---}}{---} Faithfulness Validation of Explanations}

\textbf{What we do.}
Using the counterfactual groups from Step 1 we verify that the evidence rankings
actually reflect what the model uses, via two complementary faithfulness metrics:
\begin{itemize}
  \item \textbf{Sufficiency:} Keep only the top-$k$ ranked evidence groups; mask
    everything else. If the ranking is faithful, the prediction should be largely
    preserved with just those $k$ groups. We record the absolute predicted-class
    probability and the fraction of the original probability retained.
  \item \textbf{Comprehensiveness:} Remove only the top-$k$ groups; keep everything
    else. If the ranking is faithful, the prediction probability should drop
    substantially. We record the absolute drop and fractional drop.
\end{itemize}
We sweep $k$ over $\{0,1,2,3,5,10,20\}$ and compute AUC over the resulting curves.
We run this for three rankers: counterfactual ranking (Step 1), attention-score
ranking (Step 1), and a random baseline averaged over three trials per case. Mean
and median AUC are compared across rankers.

We also run a \emph{prediction bucket analysis}: cases are split into high-confidence
correct ($\geq$0.8 confidence, right label), high-confidence wrong ($\geq$0.8,
wrong label), and low-confidence. For each bucket we report the dominant evidence
type, mean top-1 counterfactual delta, mean evidence purity (how skewed the supporting
training cases are toward one label), and mean training support size.

\textbf{Why.}
Generating importance scores is easy; demonstrating they honestly reflect the model's
computation is hard. Sufficiency and comprehensiveness are the canonical XAI
faithfulness metrics and must be reported. Comparing against attention and random
baselines determines whether our ranking extracts genuine causal information, not
artefacts of graph structure. The bucket analysis addresses what evidence the model
relies on when it is right, uncertain, and confidently wrong --- the last group is
the highest-risk failure mode in a legal setting.

\textbf{What we can infer.}
\begin{itemize}
  \item \textbf{Explanation trustworthiness:} High comprehensiveness AUC (removing
    top-$k$ hurts the model significantly) and high sufficiency AUC (keeping only
    top-$k$ preserves the prediction) together confirm the explanation captures
    genuine model reasoning.
  \item \textbf{Counterfactual vs.\ attention:} If counterfactual AUC consistently
    exceeds attention AUC, HGT attention is not a reliable importance signal despite
    its intuitive appeal --- a notable finding for the paper.
  \item \textbf{Failure mode characterisation:} The high-confidence wrong bucket
    reveals whether mistakes cluster around specific evidence types (a fixable bias)
    or are diffuse (a harder generalisation problem).
\end{itemize}

\bigskip

% -----------------------------------------------------------------------
\section{Step 3 \texorpdfstring{\textcolor{stepgray}{---}}{---} Evidence Label Skew Scoring}

\textbf{What we do.}
Using the connected-case label distribution produced in Step 1, we score every
legal authority that appears in the explanations for statistical association with
outcome. For each authority we have its training support (how many training cases
cite it) and the label breakdown (wins vs.\ losses). We run a G-test
(log-likelihood ratio test) comparing that distribution against the global training
base rate. Multiple-testing correction is applied with Benjamini-Hochberg at
$\alpha = 0.05$.

Each authority is classified as:
\begin{itemize}
  \item \textbf{Label-discriminative:} statistically significant ($q \leq 0.05$)
    and effect size $\geq$10 percentage points above or below the base rate.
  \item \textbf{Neutral:} sufficient training support but no significant skew.
  \item \textbf{Ambiguous:} insufficient training support.
\end{itemize}
We also compute log-odds of winning given the authority is present, and produce an
augmented explanations table with the skew classification attached to each evidence
row from Step 1.

\textbf{Why.}
The counterfactual experiment tells us which evidence the model relies on; this
experiment tells us whether that reliance is statistically justified by the training
data. The two are independent: a model could heavily rely on evidence with no real
label association (overfitting or structural bias) or underweight strongly
discriminative evidence (underfitting). Cross-referencing the two is essential for
interpretation.

\textbf{What we can infer.}
\begin{itemize}
  \item \textbf{Explanation grounding:} When an authority has a high counterfactual
    delta \emph{and} is label-discriminative in the direction of the prediction, we
    have strong convergent evidence of meaningful legal reasoning.
  \item \textbf{Structural vs.\ content reliance:} If the model's top-ranked
    authorities are predominantly neutral (no label association), it may be responding
    to graph topology rather than legal content.
  \item \textbf{Most predictive authorities:} The skew table ranked by log-odds is
    itself a descriptive finding about which legal authorities are most strongly
    associated with outcomes in Indian Supreme Court data.
\end{itemize}

\textbf{Outputs that feed later steps:} evidence label skew CSV (used in Steps 12
and 13 to rank distinguishing features by discriminative strength).

\bigskip

% -----------------------------------------------------------------------
\section{Step 4 \texorpdfstring{\textcolor{stepgray}{---}}{---} Case-Case Structural Community Detection}

\textbf{What we do.}
We build a TF-IDF-weighted case-feature matrix where rows are cases and columns are
legal features (statutes, provisions, precedents, judges, courts). The IDF weight of
feature $f$ is $\log(N / n_f)$ where $N$ is total cases and $n_f$ is the number of
cases containing that feature. Features appearing in fewer than 2 or more than 5,000
cases are pruned (too rare $\Rightarrow$ noise; too common $\Rightarrow$
uninformative). We compute a case-case similarity matrix as $\mathbf{W}\mathbf{W}^\top$
where $\mathbf{W}$ is the IDF-weighted feature matrix. Each case retains only its top
80 most similar neighbours and the matrix is symmetrised, giving a sparse undirected
case-case graph. Leiden community detection (resolution 1.0 by default) partitions
cases into latent legal domains. PageRank on the projection identifies representative
cases (most central in citation-overlap structure) within each community.

\textbf{Why.}
We need a principled, model-independent way to group cases before we can ask whether
the model behaves differently across legal domains. Grouping by citation structure
--- cases that cite the same laws --- clusters cases by the legal questions they
involve, independent of the model's predictions. The IDF weighting ensures
communities are shaped by specific discriminative legal content, not by ubiquitous
constitutional provisions cited in almost every case.

\textbf{What we can infer.}
\begin{itemize}
  \item \textbf{Legal domain structure:} Communities should correspond to
    interpretable legal domains. If they do, the citation network has clean latent
    structure.
  \item \textbf{Model accuracy by domain:} Domains with high accuracy and high
    confidence are well understood by the model; domains with low accuracy but high
    confidence are the most dangerous failure zones.
  \item \textbf{Train/test distribution:} The split composition of each community
    (how many train vs.\ test cases) reveals whether some legal domains are
    underrepresented at training time.
\end{itemize}

\textbf{Outputs that feed later steps:} case communities CSV with community
membership, domain bucket, and community-level accuracy (used in Steps 11, 12, 13).
Case-feature matrix (used in Steps 12, 13). Evidence label skew CSV (joint with
Step 3 from the same script).

\bigskip

% -----------------------------------------------------------------------
\section{Step 5 \texorpdfstring{\textcolor{stepgray}{---}}{---} Full Heterogeneous Graph Community Detection}

\textbf{What we do.}
Rather than projecting to a case-case graph, we run Leiden directly on the full typed
graph containing cases, statutes, provisions, and precedents as nodes simultaneously.
Because cases do not directly cite statutes (citations go through section nodes such
as arguments, facts, and preamble), we insert \emph{section bridges}: for each
``case $\to$ section $\to$ cites statute'' path we add a direct case-to-statute edge
before passing the graph to Leiden, ensuring authority nodes are connected to their
citing cases. We run Leiden at five resolution parameters (0.4, 0.7, 1.0, 1.4, 2.0)
in a single sweep, producing long- and wide-form membership tables. Edge weights are
configurable: uniform binary, or log-inverse-degree weights that downweight edges to
very high-degree hub nodes.

\textbf{Why.}
The case-case projection (Step 4) loses information about the authority nodes
themselves --- we cannot say which community a statute belongs to. Running Leiden on
the full graph gives every node a community assignment. The multi-resolution sweep
gives a hierarchy of partitions from coarse broad domains to fine sub-topics, enabling
the hierarchy analysis in Step 7.

\textbf{What we can infer.}
\begin{itemize}
  \item \textbf{Authority-case co-membership:} If a statute and the cases that cite it
    tend to land in the same community, the graph has clean domain structure. If they
    span multiple communities, that authority is a cross-domain connector (motivating
    the bridge/hub analysis in Step 8).
  \item \textbf{Resolution sensitivity:} The sweep reveals the natural scale of
    legal structure --- too low a resolution merges all of law into a few giant
    communities; too high produces noise. The elbow in the number-of-communities curve
    suggests the operating resolution.
  \item \textbf{Model error clustering:} Errors can be projected onto the full-graph
    communities, giving a richer topological view than the case-case communities
    alone.
\end{itemize}

\textbf{Outputs that feed later steps:} long-form node community membership CSV (used
in Steps 6, 7, and 8).

\bigskip

% -----------------------------------------------------------------------
\section{Step 6 \texorpdfstring{\textcolor{stepgray}{---}}{---} Full Graph Community Profiling}

\textbf{What we do.}
For a chosen resolution from the Step 5 sweep (default 1.0) we profile each community
along three axes:

\textbf{Composition:} Community size, number of test cases, counts of statutes/
provisions/precedents, dominant outcome label and legal domain bucket, full label
distribution, and the top 10 authorities per type by case count within the community.
Representative cases are identified as those whose cited authority set most overlaps
with their own community (highest same-community authority share).

\textbf{Prediction behaviour:} Overall accuracy, test-set accuracy, mean and median
confidence, and a four-bucket breakdown: (i) high-confidence correct, (ii)
high-confidence wrong (confidence $\geq$0.8, wrong label), (iii) low-confidence
correct, (iv) low-confidence wrong.

\textbf{Neighbourhood mixedness:} For every case, among all the authorities it cites,
what fraction belong to a different community? We compute the normalised entropy of
the authority community distribution and flag cases as \emph{boundary cases} if they
have enough authority citations and normalised entropy $\geq$0.5 --- cases whose
evidence spans multiple legal domains.

\textbf{Why.}
The abstract Leiden partition becomes interpretable here. Without profiling, a
community is just a set of integers; with it, it becomes ``the income tax appeals
cluster, dominated by specific provisions of the Income Tax Act, with 78\% petitioner
win rate and 82\% model accuracy''. Boundary cases are motivated by the hypothesis
that cross-domain evidence is structurally harder to reason about --- and therefore
associated with higher error rates and lower confidence.

\textbf{What we can infer.}
\begin{itemize}
  \item \textbf{Interpretable legal taxonomy:} Community profiles form a de facto
    taxonomy of Indian Supreme Court jurisprudence grounded in citation patterns ---
    independently valuable as a descriptive contribution.
  \item \textbf{Systematic failure zones:} Communities sorted by high-confidence
    wrong rate immediately flag the most dangerous failure modes.
  \item \textbf{Boundary cases as uncertainty signal:} If boundary cases tend to have
    lower model confidence and higher error rates, cross-domain evidence is a genuine
    source of model uncertainty --- a structurally grounded explanation of why some
    cases are hard.
  \item \textbf{Authority concentration:} If a small number of authorities dominate
    every community (monoculture), the dataset may generalise poorly to cases invoking
    less common legal provisions.
\end{itemize}

\bigskip

% -----------------------------------------------------------------------
\section{Step 7 \texorpdfstring{\textcolor{stepgray}{---}}{---} Community Hierarchy Analysis}

\textbf{What we do.}
For every ordered pair of resolutions $(r_{\text{coarse}}, r_{\text{fine}})$ from the
Step 5 sweep with $r_{\text{coarse}} < r_{\text{fine}}$, we compute a lineage table:
for each community at the finer resolution, which community at the coarser resolution
does it most overlap with (by node count), and what fraction of its nodes fall in that
parent? A minimum overlap threshold of 10\% suppresses noisy links. We compute
Adjusted Rand Index (ARI) and Normalised Mutual Information (NMI) between every
resolution pair for case nodes specifically. We also build full lineage chains: for
every leaf community at the finest resolution, its majority parent at every coarser
level.

\textbf{Why.}
The multi-resolution sweep produces a family of partitions but does not itself reveal
how they relate. Hierarchy analysis makes the nesting structure explicit. In an ideal
legal taxonomy, broad domains at low resolution cleanly split into sub-domains at
higher resolution. ARI and NMI between adjacent resolution levels quantify stability:
high ARI means the partition is robust to small changes in resolution parameter, which
increases confidence in the community assignments. A sharp drop identifies a structural
phase transition.

\textbf{What we can infer.}
\begin{itemize}
  \item \textbf{Hierarchical vs.\ flat structure:} High overlap between fine and
    coarse communities means the legal citation graph has genuine hierarchical
    structure (broad domains $\to$ sub-domains $\to$ specific topics).
  \item \textbf{Stable resolution range:} The range where ARI stays high between
    adjacent levels is where community assignments are trustworthy and not sensitive
    to the exact resolution chosen.
  \item \textbf{Concrete sub-domain description:} The lineage chains produce statements
    like ``community 5 at resolution 1.4 is a refinement of community 2 at resolution
    0.7 (income tax appeals), corresponding to TDS-specific provisions'' --- usable
    directly in the paper.
\end{itemize}

\bigskip

% -----------------------------------------------------------------------
\section{Step 8 \texorpdfstring{\textcolor{stepgray}{---}}{---} Bridge, Hub, and Core Authority Classification}

\textbf{What we do.}
Using the full-graph community memberships from Step 5, for every statute, provision,
and precedent we collect the Leiden community assignments of all cases that cite it
and characterise the distribution across communities using: number of distinct
communities touched, the share of citing cases in the most common community
(max-share), and the normalised entropy of the community distribution.

Three thresholds classify each authority into a role:
\begin{itemize}
  \item \textbf{Core} (max-share $\geq$80\%) --- primarily cited within one legal
    domain; a domain specialist.
  \item \textbf{Bridge} (60\%$\leq$ max-share $<$80\%, few communities touched) ---
    primarily serves one domain but also connects to one or two others.
  \item \textbf{Hub} (max-share $<$60\%, many communities) --- cited across the
    legal landscape with no dominant home domain.
  \item \textbf{Rare} ($<$5 citing cases) --- insufficient data to classify.
\end{itemize}
We also record the model's average accuracy and confidence on the cases that cite each
authority, and output a bridge-pair table showing which community pairs are most often
connected by bridge authorities.

\textbf{Why.}
From a legal perspective the taxonomy is meaningful: constitutional provisions cited
everywhere are hubs; a highly specialised tax statute is a core authority. From a
model perspective, removing a core authority should have a concentrated local effect on
its domain; removing a hub should have a diffuse global effect. If the model does not
behave this way, that tells us something important about how it processes authority
information. This classification also provides the hub authority list used in Step 10.

\textbf{What we can infer.}
\begin{itemize}
  \item \textbf{Legal network topology:} The ratio of core to bridge to hub authorities
    characterises how specialised vs.\ interconnected Indian Supreme Court jurisprudence
    is.
  \item \textbf{Authority importance conditioned on role:} Do core authorities receive
    higher counterfactual deltas within their domain than hub authorities? If yes, the
    model has implicitly learnt domain specificity; if no, it treats all legal citations
    as equally weighted regardless of domain relevance.
  \item \textbf{Cross-domain coupling:} The bridge-pair table shows which legal domain
    pairs are most tightly coupled via shared authority --- independently informative
    about Indian legal practice.
\end{itemize}

\textbf{Outputs that feed later steps:} authority role classification CSV (used in
Step 10 to select hub authorities for masking).

\bigskip

% -----------------------------------------------------------------------
\section{Step 9 \texorpdfstring{\textcolor{stepgray}{---}}{---} Identity Shortcut Audit}

\textbf{What we do.}
Independently of the model, we test whether identity names alone can predict case
outcomes from the training data. For each identity type (judge, court, petitioner,
respondent, petitioner lawyer, defence lawyer, general lawyer) we:
\begin{enumerate}
  \item Collect all training-case identity links (which training cases involve which
    judges, which petitioners, etc.).
  \item For each identity compute a smoothed posterior win rate:
    $(n_{\text{wins}} + \alpha \cdot p_{\text{prior}}) / (n_{\text{total}} + \alpha)$
    where $\alpha = 5$ is a smoothing factor and $p_{\text{prior}}$ is the global
    training win rate. This prevents single-case identities from receiving extreme
    predictions.
  \item For each test case, look up all identities it involves, weight each by its
    training support count, and compute a weighted average predicted win probability.
  \item Evaluate these identity-only predictions against ground truth using AUC-ROC,
    log-loss, Brier score, and balanced accuracy. Compare against two baselines: the
    global prior (predict the training win rate for every test case) and a domain-level
    prior (predict the domain-specific win rate), controlling for the possibility that
    identity ``signal'' is really domain signal.
  \item Run $n=100$ permutation tests (shuffling training labels within domain buckets)
    to get a p-value for whether the observed identity-only AUC exceeds chance.
\end{enumerate}
We output a top-skewed-identities table: for each type, the 50 individuals whose
training win rates deviate most from the base rate and who appear in at least 2
training and 1 test case.

\textbf{Why.}
This is a pure \emph{data} audit, separate from the model. It asks: does the dataset
contain the information needed to exploit identity shortcuts, regardless of whether the
model actually exploits it? This is necessary to correctly interpret Step 10 (mask
sensitivity). If the data contains strong identity signal (high identity-only AUC)
\emph{and} the model is sensitive to identity masking, both pieces of evidence point
to shortcut learning. If only one holds, the interpretation is more nuanced. The
smoothed-rate approach matters: identities appearing only once or twice in training
have unreliable rates and must be shrunk toward the prior.

\textbf{What we can infer.}
\begin{itemize}
  \item \textbf{Dataset-level leakage risk:} High identity-only AUC above the domain
    baseline indicates the raw dataset contains exploitable identity-outcome
    correlations, independent of legal content. This is worth reporting even if the
    model turns out not to exploit it.
  \item \textbf{Most informative identities:} The top-skewed table lists the specific
    judges, lawyers, or parties most strongly associated with outcomes. In a real legal
    system these correlations may reflect seniority, specialisation, or selection
    effects.
  \item \textbf{Cross-referencing with counterfactuals:} By joining the skewed identity
    table with the per-case counterfactual outputs (Step 1), we can ask: when the
    model assigns high importance to a judge node, is that judge one of the highly
    skewed ones? If yes, the model has specifically learnt the shortcut for high-signal
    identities.
  \item \textbf{Permutation significance:} A p-value $<0.05$ rules out the possibility
    that the observed identity-only AUC is a random artefact of the train/test split.
\end{itemize}

\bigskip

% -----------------------------------------------------------------------
\section{Step 10 \texorpdfstring{\textcolor{stepgray}{---}}{---} Mask Sensitivity Audit}

\textbf{What we do.}
We conduct two families of global post-hoc masking experiments on the frozen model,
running full-graph inference with entire categories of edges removed.

\textbf{Identity masking.} Five configurations:
\begin{itemize}
  \item Remove all edges to/from judge nodes.
  \item Remove all edges to/from party nodes (petitioner + respondent).
  \item Remove all edges to/from lawyer nodes (petitioner lawyer + defence lawyer
    + general lawyer).
  \item Remove all edges to/from court nodes.
  \item Remove all identity nodes simultaneously.
\end{itemize}
For each we record: accuracy, macro F1, mean confidence, mean and median drop in the
original predicted-class probability, and flip rate (share of cases whose predicted
class changes). Results are broken down globally and by legal domain bucket.

\textbf{Hub authority masking.} Using the authority role classification from Step 8,
we take the top 10, 25, and 50 hub authorities by citing-case count and remove all
their incident edges. We measure the same metrics.

\textbf{Why.}
Step 9 tested the data for identity signal; this step tests the \emph{model} for
identity sensitivity. This is a post-hoc \emph{inference-time} masking experiment,
not a retraining ablation: we ask what the trained model predicts when we hide judges
at test time, not what the model would learn if it never saw judges. The distinction
matters --- a model that learnt identity shortcuts during training may partially
compensate through other features at inference time, understating the problem. But the
masking approach is clean, fast, and directly comparable to the per-case counterfactual
results from Step 1.

\textbf{What we can infer.}
\begin{itemize}
  \item \textbf{Identity dependence:} Large accuracy drops and high flip rates when
    masking judges or parties are a red flag for shortcut learning. The domain-level
    breakdown reveals whether the dependence is universal or concentrated in specific
    legal areas.
  \item \textbf{Parties and lawyers as signals:} Party names and lawyer names could
    carry outcome signal (well-resourced litigants with experienced lawyers may win
    more often). Measuring the effect of masking these nodes tells us whether the model
    exploits this.
  \item \textbf{Hub authority fragility:} If accuracy collapses when the top 10 hub
    authorities are removed, the model is heavily dependent on a small number of very
    common legal references --- a robustness concern.
  \item \textbf{Convergent evidence:} Read alongside Step 9: if the data has strong
    identity signal \emph{and} the model is sensitive to masking identities, the
    evidence for shortcut learning is strong. If only one holds, the picture is more
    complex.
\end{itemize}

\bigskip

% -----------------------------------------------------------------------
\section{Step 11 \texorpdfstring{\textcolor{stepgray}{---}}{---} HGT Embedding Extraction \& Cluster Characterisation}

\textbf{What we do.}
We run a single full-graph forward pass of the frozen model on GPU (no gradient
computation) and extract the final hidden-layer representation for every case node.
These embeddings are the model's internal representation of each case after all layers
of heterogeneous message passing --- they summarise everything the model has aggregated
from the case's legal graph neighbourhood. The embeddings are L2-normalised and saved.

We then run HDBSCAN (Hierarchical Density-Based Spatial Clustering of Applications
with Noise) on the normalised embeddings. HDBSCAN is chosen over $k$-means because it
does not require specifying the number of clusters, handles clusters of varying density,
and explicitly labels low-density cases as \emph{noise} (cases the model does not
group confidently with any cluster).

We profile each embedding cluster: size, label distribution, accuracy, mean confidence,
dominant structural community (from Step 4), and domain bucket distribution. We compute
alignment metrics between the embedding clusters and the Leiden structural communities
from Step 4: ARI, NMI, homogeneity, completeness, and V-measure, both including and
excluding HDBSCAN noise points.

\textbf{Why.}
The Leiden communities (Step 4) are derived from raw citation structure; the embedding
clusters are derived from what the model has \emph{learnt} to represent. If they align,
the model has internalised legal structure. If they diverge, the model organises cases
differently --- possibly by outcome, by identity, or by some other latent factor. This
alignment test is a direct, quantitative probe of whether the model's learned
representations correspond to the structure of Indian legal citation patterns.

\textbf{What we can infer.}
\begin{itemize}
  \item \textbf{Structural faithfulness:} High ARI and NMI between embedding clusters
    and Leiden communities means the model has internalised the citation-based structure
    of the legal graph --- it groups cases the same way citation overlap does.
  \item \textbf{Alternative organisation:} Low alignment means the embedding space is
    organised by something other than citation structure. Profiling embedding clusters
    by label distribution, confidence, and domain immediately suggests hypotheses
    (e.g.\ the model has collapsed to a simple win/loss geometry, or groups by which
    court presided).
  \item \textbf{Noise cases:} HDBSCAN noise points are cases in a low-density region
    of embedding space --- not similar to any dense group. These may be the hardest
    cases for the model, or the most legally unusual. Their accuracy and confidence
    profiles are worth examining separately.
  \item \textbf{Community--cluster mismatch:} If a structural community maps to many
    distinct embedding clusters, the model makes further distinctions within that legal
    domain that the citation graph alone does not capture --- potentially meaningful
    or potentially indicative of spurious features.
\end{itemize}

\textbf{Outputs that feed later steps:} normalised case embeddings (used in Steps
12 and 13).

\bigskip

% -----------------------------------------------------------------------
\section{Step 12 \texorpdfstring{\textcolor{stepgray}{---}}{---} Counterfactual Neighborhoods in Embedding Space}

\textbf{What we do.}
For every test case we find its nearest \emph{opposite-label} training case in the
normalised embedding space using GPU-accelerated batched cosine similarity. The search
is label-aware: for a test case labelled WIN we search among training LOSS cases, and
vice versa. This gives, for each test case, the training case most similar in
embedding space but with the opposite ground-truth label --- the case that most nearly
``confused'' the model at training time.

Having found the nearest opposite-label neighbour, we compare the two cases' evidence
feature sets from the TF-IDF feature matrix (Step 4) and identify:
\begin{itemize}
  \item Features present in the test case but absent in the neighbour (\emph{query-only}).
  \item Features present in the neighbour but absent in the test case (\emph{opposite-only}).
  \item Features shared by both.
\end{itemize}
For each difference feature we look up its label skew from Step 3 and rank the
distinguishing features by a combination of discriminative strength (log-odds) and
rarity (IDF weight). The top distinguishing features, the cosine similarity to the
neighbour, and a narrative summary are written per case.

\textbf{Why.}
The typed counterfactual masking (Step 1) explains a prediction in absolute terms:
``this statute was important for this case''. Counterfactual neighbourhood analysis
asks a relative question: ``what is different between this case and the most similar
case with the opposite outcome?'' This directly probes the model's decision boundary
rather than its interior. Cases with a nearby opposite-label neighbour are close to
the margin in the model's representation space --- the features that distinguish them
reveal what the model is using to separate WIN from LOSS at the boundary.

\textbf{What we can infer.}
\begin{itemize}
  \item \textbf{Decision boundary characterisation:} If the distinguishing features
    are predominantly label-discriminative legal authorities (Step 3), the model's
    decision boundary is aligned with meaningful legal content. If they are identity
    nodes (different judge, same statutes), the boundary is shaped by shortcuts.
  \item \textbf{Cosine similarity as a margin signal:} Test cases with very high
    cosine similarity ($\geq$0.95) to their nearest opposite-label neighbour are the
    model's hardest cases: nearly identical internal representations, opposite labels.
    These deserve qualitative examination.
  \item \textbf{Systematic vs.\ idiosyncratic margins:} If the same distinguishing
    features recur across many test cases' nearest-neighbour comparisons, those features
    systematically shape the global decision boundary. If every case has unique
    distinguishing features, the margins are idiosyncratic.
  \item \textbf{Community membership of pairs:} If a test case and its nearest
    opposite-label neighbour are in the same structural community (similar citations,
    Step 4) but have opposite labels and similar embeddings, the model has not found
    a way to separate nearly-identical cases with different outcomes --- possibly
    because the legal question is genuinely ambiguous.
\end{itemize}

\bigskip

% -----------------------------------------------------------------------
\section{Step 13 \texorpdfstring{\textcolor{stepgray}{---}}{---} Per-Case Traceability Reports}

\textbf{What we do.}
As the final synthesis step, for any selected set of cases we generate self-contained
HTML reports that bring together the outputs of all previous steps into a
human-readable form. Each report contains:
\begin{itemize}
  \item Header with case identifier, split, target label, predicted label, and
    confidence.
  \item A \emph{structured reading} paragraph synthesised automatically from the top
    evidence group: which evidence was most important, the direction of its effect,
    whether it flips the prediction, and what the connected training cases look like.
  \item An interactive force-directed graph (rendered in-browser) showing the typed
    computation subgraph --- the target case at the centre, evidence nodes sized by
    counterfactual delta (Step 1), edges coloured by relation type and weighted by
    attention score (Step 1).
  \item An evidence group table with counterfactual rank (Step 1), attention rank
    (Step 1), a ``both agree'' marker when the two rankings concur, delta values,
    attention scores, and the connected training-case label distribution (Step 1).
    Label skew class from Step 3 is attached to each row.
  \item A provenance trail: for every top-ranked evidence node, a list of training
    cases that share it (up to a configurable limit), with links to the source JSON
    judgment files, the training-case label, and the typed graph path connecting the
    training case to the evidence node.
  \item Machine-readable JSON and DOT graph exports alongside the HTML.
\end{itemize}
Case selection can be targeted (specific indices), or ranked by counterfactual
importance so the most ``interesting'' cases are always surfaced first.

\textbf{Why.}
XAI research requires both quantitative faithfulness metrics and human-interpretable
outputs. Metrics prove the system is not lying; human-readable reports make it usable.
The traceability reports bridge this gap. They are designed so that a lawyer, a
supervisor, or a judge reviewing the model's reasoning can follow every claim back to
its source --- the specific training cases that established the association between
evidence and outcome. The provenance trail is the difference between ``the model
relied on this statute because it is associated with wins'' and ``the model relied on
this statute because cases A, B, and C all cited it and were won by the petitioner''
--- the latter is auditable and falsifiable.

\textbf{What we can infer.}
\begin{itemize}
  \item \textbf{Face validity:} Do the explanations make legal sense? A qualitative
    review of reports on a representative sample is an important sanity check that
    quantitative metrics alone cannot provide.
  \item \textbf{Attention vs.\ counterfactual disagreements:} When the ``both agree''
    marker is absent, the report makes the disagreement visible case-by-case, enabling
    inspection of what structural feature drives the discrepancy.
  \item \textbf{Provenance quality:} If the model highlights a statute and the
    provenance trail shows it was cited in training cases with a 90\% win rate, that is
    meaningful legal grounding. If training cases are mixed, the model's reliance is
    less interpretable.
  \item \textbf{Error diagnosis at the instance level:} Reports for high-confidence
    wrong cases (flagged in Step 2) let us ask: what did the model see in these cases,
    and why was that evidence insufficient or misleading? This closes the loop between
    the quantitative failure analysis (Steps 1--2) and concrete case-level inspection.
\end{itemize}

\end{document}
